\IfFileExists{stacks-project.cls}{%
\documentclass{stacks-project}
}{%
\documentclass{amsart}
}
\usepackage{amssymb}

% For dealing with references we use the comment environment
\usepackage{verbatim}
\usepackage{fancyvrb}
\newenvironment{reference}{\comment}{\endcomment}
%\newenvironment{reference}{}{}
\newenvironment{slogan}{\comment}{\endcomment}
\newenvironment{history}{\comment}{\endcomment}

% For source figures retained as assets
\usepackage{graphicx}

% For commutative diagrams we use Xy-pic
\usepackage[all]{xy}

% We use 2cell for 2-commutative diagrams.
\xyoption{2cell}
\UseAllTwocells

% We use multicol for the list of chapters between chapters
\usepackage{multicol}

% This is generally recommended for better output
\usepackage{lmodern}
\usepackage[T1]{fontenc}
\usepackage{textalpha}

% For cross-file-references
% We prefer the snapshot tmp/xref/#2.aux,
% which latexmkrc replaces atomically, so
% that documents compiled simultaneously
% never read a half-written .aux file.
% The last argument keeps links to #2.pdf.
\usepackage{xr-hyper}
\newcommand{\maybeexternaldocument}[2]{%
  \IfFileExists{#2.aux}{%
    \externaldocument[#1]{#2}%
  }{%
    \PackageWarning{stack}{Missing #2.aux; external references with prefix #1 unavailable}%
  }%
}
\let\externaldocumentorig\externaldocument
\renewcommand{\externaldocument}[2][]{%
  \IfFileExists{tmp/xref/#2.aux}{%
    \externaldocumentorig[#1]{tmp/xref/#2}[#2.pdf]%
  }{%
  \IfFileExists{#2.aux}{%
    \externaldocumentorig[#1]{#2}%
  }{%
    \PackageWarning{stack}{Missing #2.aux; external references with prefix #1 unavailable}%
  }%
  }%
}

% Package for hypertext links:
\usepackage{hyperref}

% For any local file, say "hello.tex" you want to link to please
% use \externaldocument[hello-]{hello}
\externaldocument[introduction-]{introduction}
\externaldocument[conventions-]{conventions}
\externaldocument[sets-]{sets}
\externaldocument[categories-]{categories}
\externaldocument[topology-]{topology}
\externaldocument[sheaves-]{sheaves}
\externaldocument[sites-]{sites}
\externaldocument[stacks-]{stacks}
\externaldocument[fields-]{fields}
\externaldocument[algebra-]{algebra}
\externaldocument[brauer-]{brauer}
\externaldocument[homology-]{homology}
\externaldocument[derived-]{derived}
\externaldocument[simplicial-]{simplicial}
\externaldocument[more-algebra-]{more-algebra}
\externaldocument[smoothing-]{smoothing}
\externaldocument[modules-]{modules}
\externaldocument[sites-modules-]{sites-modules}
\externaldocument[injectives-]{injectives}
\externaldocument[cohomology-]{cohomology}
\externaldocument[sites-cohomology-]{sites-cohomology}
\externaldocument[dga-]{dga}
\externaldocument[dpa-]{dpa}
\externaldocument[sdga-]{sdga}
\externaldocument[hypercovering-]{hypercovering}
\externaldocument[schemes-]{schemes}
\externaldocument[constructions-]{constructions}
\externaldocument[properties-]{properties}
\externaldocument[morphisms-]{morphisms}
\externaldocument[coherent-]{coherent}
\externaldocument[divisors-]{divisors}
\externaldocument[limits-]{limits}
\externaldocument[varieties-]{varieties}
\externaldocument[topologies-]{topologies}
\externaldocument[descent-]{descent}
\externaldocument[perfect-]{perfect}
\externaldocument[more-morphisms-]{more-morphisms}
\externaldocument[flat-]{flat}
\externaldocument[groupoids-]{groupoids}
\externaldocument[more-groupoids-]{more-groupoids}
\externaldocument[etale-]{etale}
\externaldocument[chow-]{chow}
\externaldocument[intersection-]{intersection}
\externaldocument[pic-]{pic}
\externaldocument[weil-]{weil}
\externaldocument[adequate-]{adequate}
\externaldocument[dualizing-]{dualizing}
\externaldocument[duality-]{duality}
\externaldocument[discriminant-]{discriminant}
\externaldocument[derham-]{derham}
\externaldocument[local-cohomology-]{local-cohomology}
\externaldocument[algebraization-]{algebraization}
\externaldocument[curves-]{curves}
\externaldocument[resolve-]{resolve}
\externaldocument[models-]{models}
\externaldocument[functors-]{functors}
\externaldocument[equiv-]{equiv}
\externaldocument[pione-]{pione}
\externaldocument[etale-cohomology-]{etale-cohomology}
\externaldocument[proetale-]{proetale}
\externaldocument[relative-cycles-]{relative-cycles}
\externaldocument[more-etale-]{more-etale}
\externaldocument[trace-]{trace}
\externaldocument[crystalline-]{crystalline}
\externaldocument[spaces-]{spaces}
\externaldocument[spaces-properties-]{spaces-properties}
\externaldocument[spaces-morphisms-]{spaces-morphisms}
\externaldocument[decent-spaces-]{decent-spaces}
\externaldocument[spaces-cohomology-]{spaces-cohomology}
\externaldocument[spaces-limits-]{spaces-limits}
\externaldocument[spaces-divisors-]{spaces-divisors}
\externaldocument[spaces-over-fields-]{spaces-over-fields}
\externaldocument[spaces-topologies-]{spaces-topologies}
\externaldocument[spaces-descent-]{spaces-descent}
\externaldocument[spaces-perfect-]{spaces-perfect}
\externaldocument[spaces-more-morphisms-]{spaces-more-morphisms}
\externaldocument[spaces-flat-]{spaces-flat}
\externaldocument[spaces-groupoids-]{spaces-groupoids}
\externaldocument[spaces-more-groupoids-]{spaces-more-groupoids}
\externaldocument[bootstrap-]{bootstrap}
\externaldocument[spaces-pushouts-]{spaces-pushouts}
\externaldocument[spaces-chow-]{spaces-chow}
\externaldocument[groupoids-quotients-]{groupoids-quotients}
\externaldocument[spaces-more-cohomology-]{spaces-more-cohomology}
\externaldocument[spaces-simplicial-]{spaces-simplicial}
\externaldocument[spaces-duality-]{spaces-duality}
\externaldocument[formal-spaces-]{formal-spaces}
\externaldocument[restricted-]{restricted}
\externaldocument[spaces-resolve-]{spaces-resolve}
\externaldocument[formal-defos-]{formal-defos}
\externaldocument[defos-]{defos}
\externaldocument[cotangent-]{cotangent}
\externaldocument[examples-defos-]{examples-defos}
\externaldocument[algebraic-]{algebraic}
\externaldocument[examples-stacks-]{examples-stacks}
\externaldocument[stacks-sheaves-]{stacks-sheaves}
\externaldocument[criteria-]{criteria}
\externaldocument[artin-]{artin}
\externaldocument[quot-]{quot}
\externaldocument[stacks-properties-]{stacks-properties}
\externaldocument[stacks-morphisms-]{stacks-morphisms}
\externaldocument[stacks-limits-]{stacks-limits}
\externaldocument[stacks-cohomology-]{stacks-cohomology}
\externaldocument[stacks-perfect-]{stacks-perfect}
\externaldocument[stacks-introduction-]{stacks-introduction}
\externaldocument[stacks-more-morphisms-]{stacks-more-morphisms}
\externaldocument[stacks-geometry-]{stacks-geometry}
\externaldocument[moduli-]{moduli}
\externaldocument[moduli-curves-]{moduli-curves}
\externaldocument[examples-]{examples}
\externaldocument[exercises-]{exercises}
\externaldocument[guide-]{guide}
\externaldocument[desirables-]{desirables}
\externaldocument[coding-]{coding}
\externaldocument[obsolete-]{obsolete}
\externaldocument[fdl-]{fdl}
\externaldocument[index-]{index}

% Heap Project documents
\externaldocument[k3-]{k3}
\externaldocument[torelli-]{torelli}
\externaldocument[linear-algebra-]{linear-algebra}
\externaldocument[groups-]{groups}
\externaldocument[complex-analysis-]{complex-analysis}
\externaldocument[functional-analysis-]{functional-analysis}
\externaldocument[categories-]{categories}
\externaldocument[commutative-algebra-]{commutative-algebra}
\externaldocument[ringed-spaces-]{ringed-spaces}
\externaldocument[homological-algebra-]{homological-algebra}
\externaldocument[lie-algebras-]{lie-algebras}
\externaldocument[representation-theory-]{representation-theory}
\externaldocument[smooth-manifolds-]{smooth-manifolds}
\externaldocument[differential-topology-]{differential-topology}
\externaldocument[algebraic-topology-]{algebraic-topology}
\externaldocument[differential-geometry-]{differential-geometry}
\externaldocument[algebraic-geometry-]{algebraic-geometry}
\externaldocument[symplectic-geometry-]{symplectic-geometry}
\externaldocument[complex-geometry-]{complex-geometry}
\externaldocument[hodge-theory-]{hodge-theory}
\externaldocument[%
variations-of-hodge-structures-]{%
variations-of-hodge-structures}
\externaldocument[%
homological-mirror-symmetry-]{%
homological-mirror-symmetry}
\maybeexternaldocument{deformations-}{deformations}
\externaldocument[vertex-algebras-]{vertex-algebras}
\externaldocument[springer-]{springer}
\externaldocument[infty-categories-]{infty-categories}
\externaldocument[seminars-differential-geometry-]{%
seminars-differential-geometry}
\externaldocument[seminars-complex-geometry-]{%
seminars-complex-geometry}
\externaldocument[seminars-algebraic-geometry-]{%
seminars-algebraic-geometry}
\externaldocument[seminars-others-]{%
seminars-others}
\externaldocument[%
geometric-prequantization-]{%
geometric-prequantization}
\externaldocument[surfaces-]{surfaces}
\externaldocument[birational-maps-commutative-algebra-]{birational-maps-commutative-algebra}
\externaldocument[cremona-transformations-]{cremona-transformations}
\externaldocument[derived-categories-of-sheaves-]{derived-categories-of-sheaves}
\externaldocument[geometric-invariant-theory-]{geometric-invariant-theory}
\externaldocument[geometric-stability-conditions-and-group-actions-]{geometric-stability-conditions-and-group-actions}
\externaldocument[moduli-spaces-of-sheaves-]{moduli-spaces-of-sheaves}
\externaldocument[bridgeland-stability-general-theory-]{bridgeland-stability-general-theory}
\externaldocument[hilbert-schemes-]{hilbert-schemes}
\externaldocument[noncommutative-abelian-surfaces-and-kummer-type-hyperkahler-varieties-]{noncommutative-abelian-surfaces-and-kummer-type-hyperkahler-varieties}
\externaldocument[mumford-tate-groups-in-hodge-theory-]{mumford-tate-groups-in-hodge-theory}
\externaldocument[hodge-birational-atoms-2026-]{hodge-birational-atoms-2026}
\externaldocument[xxii-egd-2026-]{%
xxii-egd-2026}
\externaldocument[pde-]{pde}
\externaldocument[probability-]{probability}

\externaldocument[linguistics-]{linguistics}
\externaldocument[genealogia-familiar-]{%
genealogia-familiar}

\externaldocument[%
16th-alga-meeting-2026-]{%
16th-alga-meeting-2026}

\externaldocument[%
iv-jornadas-lefschetz-2026-]{%
iv-jornadas-lefschetz-2026}

\externaldocument[reading-the-masters-]{%
reading-the-masters}

\externaldocument[real-analysis-]{%
real-analysis}

\externaldocument[optimal-transport-]{%
optimal-transport}

\externaldocument[history-]{history}

\externaldocument[seiberg-witten-]{seiberg-witten}

\externaldocument[history-philosophy-biology-]{history-philosophy-biology}

\externaldocument[chemistry-]{chemistry}

\externaldocument[molecular-cell-biology-]{molecular-cell-biology}

\externaldocument[microbiology-]{microbiology}

\externaldocument[protistology-]{protistology}

\externaldocument[mycology-]{mycology}

\externaldocument[botany-]{botany}

\externaldocument[zoology-]{zoology}

\externaldocument[animal-histology-]{animal-histology}

\externaldocument[animal-physiology-]{animal-physiology}

\externaldocument[systematics-biogeography-]{systematics-biogeography}

\externaldocument[ecology-natural-resources-]{ecology-natural-resources}

\externaldocument[earth-sciences-]{earth-sciences}

\externaldocument[biostatistics-]{biostatistics}

\externaldocument[mathematical-systems-biology-]{mathematical-systems-biology}

\externaldocument[cimpa-el-salvador-2024-]{cimpa-el-salvador-2024}

% Theorem environments.
%
\theoremstyle{plain}
\newtheorem{theorem}[subsection]{Theorem}
\newtheorem{proposition}[subsection]{Proposition}
\newtheorem{lemma}[subsection]{Lemma}
\newtheorem{corollary}[subsection]{Corollary}

\theoremstyle{definition}
\newtheorem{definition}[subsection]{Definition}
\newtheorem{example}[subsection]{Example}
\newtheorem{exercise}[subsection]{Exercise}
\newtheorem{situation}[subsection]{Situation}

\theoremstyle{remark}
\newtheorem{remark}[subsection]{Remark}
\newtheorem{remarks}[subsection]{Remarks}

\numberwithin{equation}{subsection}

% Macros
%
\def\lim{\mathop{\mathrm{lim}}\nolimits}
\def\colim{\mathop{\mathrm{colim}}\nolimits}
\def\Spec{\mathop{\mathrm{Spec}}}
\def\Hom{\mathop{\mathrm{Hom}}\nolimits}
\def\Ext{\mathop{\mathrm{Ext}}\nolimits}
\def\SheafHom{\mathop{\mathcal{H}\!\mathit{om}}\nolimits}
\def\SheafExt{\mathop{\mathcal{E}\!\mathit{xt}}\nolimits}
\def\Sch{\mathit{Sch}}
\def\Mor{\mathop{\mathrm{Mor}}\nolimits}
\def\Ob{\mathop{\mathrm{Ob}}\nolimits}
\def\Sh{\mathop{\mathit{Sh}}\nolimits}
\def\NL{\mathop{N\!L}\nolimits}
\def\CH{\mathop{\mathrm{CH}}\nolimits}
\def\proetale{{pro\text{-}\acute{e}tale}}
\def\etale{{\acute{e}tale}}
\def\QCoh{\mathit{QCoh}}
\def\Ker{\mathop{\mathrm{Ker}}}
\def\Im{\mathop{\mathrm{Im}}}
\def\Coker{\mathop{\mathrm{Coker}}}
\def\Coim{\mathop{\mathrm{Coim}}}

% Boxtimes
%
\DeclareMathSymbol{\boxtimes}{\mathbin}{AMSa}{"02}

%
% Macros for moduli stacks/spaces
%
\def\QCohstack{\mathcal{QC}\!\mathit{oh}}
\def\Cohstack{\mathcal{C}\!\mathit{oh}}
\def\Spacesstack{\mathcal{S}\!\mathit{paces}}
\def\Quotfunctor{\mathrm{Quot}}
\def\Hilbfunctor{\mathrm{Hilb}}
\def\Curvesstack{\mathcal{C}\!\mathit{urves}}
\def\Polarizedstack{\mathcal{P}\!\mathit{olarized}}
\def\Complexesstack{\mathcal{C}\!\mathit{omplexes}}
% \Pic is the operator that assigns to X its picard group, usage \Pic(X)
% \Picardstack_{X/B} denotes the Picard stack of X over B
% \Picardfunctor_{X/B} denotes the Picard functor of X over B
\def\Pic{\mathop{\mathrm{Pic}}\nolimits}
\def\Picardstack{\mathcal{P}\!\mathit{ic}}
\def\Picardfunctor{\mathrm{Pic}}
\def\Deformationcategory{\mathcal{D}\!\mathit{ef}}


\begin{document}

\title{Animal Physiology}
\maketitle

\phantomsection
\label{section-phantom}

\tableofcontents

\section{What this chapter covers}
\label{section-overview}

These are notes on comparative animal
physiology: regulation, excitable
membranes, synapses, the senses across
animals, the nervous and endocrine
systems, and biological clocks. A short
section on water, salts and excretion
closes the chapter.

\textbf{Sources of these notes.}
The notes rest on the course
\emph{Fisiolog\'ia Animal} (listed on the
homework as \emph{Biolog\'ia de Animales
II}), fifth semester, Fall 2015
(semester 2016-1), Facultad de Ciencias,
UNAM. What survives is Daniel's answered
homework for two blocks (regulation,
membranes and synapses; senses, nervous
system, hormones and clocks), the
teacher's question sheets, a lab report
on a cricket electroretinogram, a short
answer on external potassium, a
spreadsheet of membrane potentials, the
data and statistical analyses of three
practicals, his notes of a lab class, and
journal papers read in class. No textbook
is cited in the coursework. The notes
were checked and completed with Hill,
Wyse and Anderson
\cite{hill2012animal}, the Eckert text by
Randall, Burggren and French
\cite{randall1997eckert}, Willmer, Stone
and Johnston \cite{willmer2005environmental},
Schmidt's sensory physiology
\cite{schmidt1981sensory} and Mrosovsky
\cite{mrosovsky1990rheostasis}.

\section{Homeostasis and rheostasis}
\label{section-homeostasis}

Claude Bernard noticed that mammals keep
their blood and tissue fluids nearly
constant (glucose, temperature, ions,
pH), and Walter Cannon coined the word
\emph{homeostasis} for this constancy
together with the regulatory systems that
restore it after a disturbance
\cite[Chapter 1]{hill2012animal}.

In a control system, the \emph{controlled
variable} is the property being kept
nearly constant and the \emph{set point}
is the level at which it is kept. There
is \emph{feedback} when the system uses
information on the controlled variable
itself to govern its actions.
\begin{itemize}
\item In \emph{negative feedback} the
system opposes deviations from the set
point (blood glucose: the liver takes up
or releases glucose). Negative feedback
is virtually synonymous with homeostasis.
\item In \emph{positive feedback} the
system reinforces the deviation until
something breaks the cycle: the rising
phase of the action potential, or
uterine contractions and hormone release
during birth.
\end{itemize}
Norbert Wiener's \emph{cybernetics} is the
general theory of such self-regulating
systems, in machines and in organisms
alike.

Not every animal regulates every
variable: a salmon is a temperature
conformer but a chloride regulator.
Chronic responses to a changed
environment are called acclimation when
the new environment differs in a few
defined ways (a laboratory phenomenon)
and acclimatization when it is a
different natural environment, such as
winter versus summer or a high altitude
\cite[Chapter 1]{hill2012animal}.

\begin{definition}
\label{definition-rheostasis}
\emph{Rheostasis} is a condition in which,
at any one instant, homeostatic defenses
are present, but over a span of time the
level that is defended changes. It is
\emph{programmed} when the change is
obligatory at certain phases of the life
cycle, and \emph{reactive} when it occurs
in response to stimuli that may or may
not be encountered.
\end{definition}

This is Mrosovsky's definition
\cite[Chapter 3]{mrosovsky1990rheostasis}.
Programmed examples are the daily cycle
of body temperature and the resetting of
the gonadostat over the estrous cycle
\cite[Chapter 4]{mrosovsky1990rheostasis};
the classic reactive example is fever,
in which a higher body temperature is
actively defended
\cite[Chapter 5]{mrosovsky1990rheostasis}.
Mrosovsky stresses that the two kinds
form a continuum.

\section{Membrane potentials}
\label{section-membrane-potential}

All cells keep more potassium and less
sodium and chloride inside than outside,
because the Na$^+$-K$^+$-ATPase pumps
Na$^+$ out and K$^+$ in; neither ion is
at equilibrium \cite[Chapter
12]{hill2012animal}.

\textbf{Nernst equation.}
For a single ion species $X$ of valence
$z$ that can cross the membrane, the
membrane potential (inside minus
outside) at which $X$ is at equilibrium
is
\begin{equation}
\label{equation-nernst}
E_X = \frac{RT}{zF}
\ln\frac{[X]_o}{[X]_i},
\end{equation}
with $R$ the gas constant, $T$ the
absolute temperature, $F$ Faraday's
constant and $[X]_o$, $[X]_i$ the outside
and inside concentrations. For a
monovalent ion $RT/F$ is about
$25.3\ \mathrm{mV}$ at
$20\,^{\circ}\mathrm{C}$,
so that
$$
E_X \approx \frac{58\ \mathrm{mV}}{z}
\log_{10}\frac{[X]_o}{[X]_i}
$$
near room temperature, and $61\
\mathrm{mV}$ replaces $58\ \mathrm{mV}$
at mammalian body temperature
\cite[Chapter 12]{hill2012animal},
\cite[Chapter 5]{randall1997eckert}.

\textbf{Goldman-Hodgkin-Katz equation.}
If several ions permeate, each
contributes in proportion to its
permeability $P$. For K$^+$, Na$^+$ and
Cl$^-$ the steady membrane potential is
\begin{equation}
\label{equation-ghk}
V_m = \frac{RT}{F}\ln
\frac{P_{\mathrm{K}}[\mathrm{K}^+]_o
+ P_{\mathrm{Na}}[\mathrm{Na}^+]_o
+ P_{\mathrm{Cl}}[\mathrm{Cl}^-]_i}
{P_{\mathrm{K}}[\mathrm{K}^+]_i
+ P_{\mathrm{Na}}[\mathrm{Na}^+]_i
+ P_{\mathrm{Cl}}[\mathrm{Cl}^-]_o}.
\end{equation}
The chloride concentrations appear the
other way round because the charge of
chloride is negative. With only one
permeant ion the equation reduces to the
Nernst equation for that ion. At rest
$P_{\mathrm{K}}$ dominates, so $V_m$ lies
near $E_{\mathrm{K}}$ \cite[Chapter
12]{hill2012animal}, \cite[Chapter
5]{randall1997eckert}.

\textbf{Parallel-conductance model.}
Hodgkin and Huxley represented the
membrane as a capacitor $C_m$ in parallel
with one branch per ion: a conductance
$g_X$ (open channels) in series with a
battery $E_X$. Each ionic current obeys
Ohm's law with the driving force
$V_m - E_X$:
\begin{equation}
\label{equation-ohm-ion}
I_X = g_X\,(V_m - E_X)
\end{equation}
\cite[Chapter 5]{randall1997eckert}. At
rest $V_m$ is constant, so no current
charges the capacitor and the ionic
currents add up to zero. Solving
$\sum_X g_X (V_m - E_X) = 0$ gives
\begin{equation}
\label{equation-parallel-conductance}
V_m = \frac{\sum_X g_X E_X}{\sum_X g_X},
\end{equation}
a weighted mean of the equilibrium
potentials, weighted by conductances.
\emph{Conductance} is the reciprocal of
resistance, measured in siemens (the old
unit mho is the same thing).

\textbf{External potassium.}
Raising $[\mathrm{K}^+]_o$ makes
$E_{\mathrm{K}}$ less negative but leaves
$E_{\mathrm{Na}}$ and $E_{\mathrm{Cl}}$
unchanged, since each equilibrium
potential depends only on that ion's own
gradient; $V_m$ still depolarizes because
it depends on all ions. Daniel's
spreadsheet for this exercise shows it:
as $[\mathrm{K}^+]_o$ rises from $5$ to
$116\ \mathrm{mM}$, $E_{\mathrm{K}}$ goes
from $-85$ to $-4\ \mathrm{mV}$ and $V_m$
from $-67$ to $-11\ \mathrm{mV}$, while
$E_{\mathrm{Na}} = +54\ \mathrm{mV}$ and
$E_{\mathrm{Cl}} = -75\ \mathrm{mV}$ stay
fixed.

\section{Example: the Aplysia giant cell}
\label{section-aplysia}

\begin{example}
\label{example-aplysia}
The homework gives values reported by
Brown and coworkers (1970) for the giant
neuron of the sea hare \emph{Aplysia}:
\begin{center}
\begin{tabular}{lccc}
Ion & in (mM) & out (mM)
& $g$ ($\mu$S) \\
\hline
Na$^+$ & 50 & 337 & 0.11 \\
K$^+$ & 168 & 6 & 0.57 \\
Cl$^-$ & 41 & 340 & 0.32 \\
\end{tabular}
\end{center}
At $20\,^{\circ}\mathrm{C}$,
$T = 293.15\ \mathrm{K}$ and $RT/F =
25.26\ \mathrm{mV}$. The Nernst equation
\eqref{equation-nernst} gives
$$
\begin{aligned}
E_{\mathrm{Na}} &= 25.26
\ln\frac{337}{50}
= +48.2\ \mathrm{mV}, \\
E_{\mathrm{K}} &= 25.26
\ln\frac{6}{168}
= -84.2\ \mathrm{mV}, \\
E_{\mathrm{Cl}} &= -25.26
\ln\frac{340}{41}
= -53.4\ \mathrm{mV}.
\end{aligned}
$$
The conductances add up to $1.00\
\mu\mathrm{S}$, so by
\eqref{equation-parallel-conductance}
$$
\begin{aligned}
V_m &= 0.11(48.2) + 0.57(-84.2)
+ 0.32(-53.4) \\
&= 5.3 - 48.0 - 17.1
\approx -60\ \mathrm{mV}.
\end{aligned}
$$
The resting potential lies between
$E_{\mathrm{K}}$ and $E_{\mathrm{Cl}}$,
pulled up a little by sodium.
\end{example}

\section{Cable properties and conduction}
\label{section-cable}

\textbf{Time constant.}
Inject a step of current into a cell. The
current first redistributes charge on the
membrane capacitance (capacitive current)
and then flows through the channels
(resistive current), so the voltage
change approaches its final value
$\Delta V_\infty$ exponentially:
$$
\Delta V(t) = \Delta V_\infty
\left(1 - e^{-t/\tau}\right),
\qquad \tau = R_m C_m .
$$
The \emph{time constant} $\tau$ is the
time to reach $63\%$ of $\Delta V_\infty$;
for many cells it is $2$ to $20\
\mathrm{ms}$. It is measured from the
recorded voltage curve
\cite[Chapter 12]{hill2012animal}.

\textbf{Length constant.}
Along an axon, some current leaks out
through the membrane, so a steady voltage
change decays with distance $x$ from the
source (electrotonic conduction):
$$
\Delta V(x) = \Delta V_0\, e^{-x/\lambda},
\qquad
\lambda = \sqrt{\frac{r_m}{r_i + r_o}},
$$
with $r_m$ the membrane resistance of a
unit length of axon and $r_i$, $r_o$ the
internal and external longitudinal
resistances per unit length. The
\emph{length constant} $\lambda$ is the
distance at which $\Delta V$ has fallen to
$37\%$ of $\Delta V_0$; it ranges from
about $0.1$ to $5\ \mathrm{mm}$
\cite[Chapter 6]{randall1997eckert}. These
are the \emph{cable properties} of the
neuron.

\textbf{Action potential.}
A depolarization to threshold opens
voltage-gated Na$^+$ channels; Na$^+$
entry depolarizes further and opens more
channels (positive feedback), driving
$V_m$ toward $E_{\mathrm{Na}}$. Sodium
channels then inactivate and
voltage-gated K$^+$ channels open, which
repolarizes the membrane and briefly
hyperpolarizes it toward $E_{\mathrm{K}}$
\cite[Chapter 12]{hill2012animal},
\cite[Chapter 5]{randall1997eckert}.

\textbf{Conduction velocity.}
The inward current at an active patch
spreads electrotonically and brings the
next patch to threshold. Anything that
increases $\lambda$ or shortens the time to
threshold speeds conduction. The three
main factors are:
\begin{itemize}
\item axon diameter: $r_m/r_i$ grows
linearly with diameter $D$, so in
unmyelinated axons the velocity grows
roughly like $D^{1/2}$ (hence the giant
axons of squid and other animals);
\item myelination: the insulating sheath
raises membrane resistance, and the
action potential jumps between nodes of
Ranvier (saltatory conduction);
\item temperature: warmer axons conduct
faster.
\end{itemize}
See \cite[Chapter 12]{hill2012animal}.

\section{Synapses}
\label{section-synapses}

\textbf{Electrical versus chemical.}
At an electrical synapse, current passes
directly between cells through gap
junctions (channels built from connexins
in vertebrates), so transmission is
almost instantaneous. At a chemical
synapse the cells are separated by a
synaptic cleft of $20$ to $30\
\mathrm{nm}$; the presynaptic terminal
releases neurotransmitter in quantal
packets, which acts on postsynaptic
receptors, after a short synaptic delay
\cite[Chapter 13]{hill2012animal}.

\textbf{Ionotropic versus metabotropic.}
An ionotropic receptor is a ligand-gated
channel: the same molecule binds the
transmitter and conducts ions, giving
fast synaptic potentials. A metabotropic
receptor is a G protein-coupled receptor
with seven transmembrane segments; it
acts through second messengers (for
example cAMP via adenylyl cyclase), more
slowly and for longer, and can change
excitability, metabolism or gene
expression \cite[Chapter
13]{hill2012animal}.

\textbf{Main transmitters.}
The table follows
\cite[Table 13.2]{hill2012animal}.
\begin{center}
\begin{tabular}{lll}
Transmitter & Receptor & Type \\
\hline
acetylcholine & nicotinic & ionotr. \\
 & muscarinic M1-M5 & metabotr. \\
norepinephrine & $\alpha$, $\beta$
& metabotr. \\
dopamine & D1-D3 & metabotr. \\
serotonin & 5-HT1, 5-HT2
& metabotr. \\
 & 5-HT3 & ionotr. \\
glutamate & AMPA, NMDA & ionotr. \\
 & mGluR & metabotr. \\
GABA & GABA$_{\mathrm{A}}$ & ionotr. \\
 & GABA$_{\mathrm{B}}$ & metabotr. \\
glycine & glycine receptor
& ionotr. \\
\end{tabular}
\end{center}
Nicotinic receptors mediate the
vertebrate neuromuscular junction;
glutamate is the main fast excitatory and
GABA and glycine the main fast inhibitory
transmitters. Neuropeptides, such as
substance P and the enkephalins, act on
metabotropic receptors and mediate slow
effects \cite[Chapter
6]{randall1997eckert}, \cite[Chapter
13]{hill2012animal}.

\section{Sensory physiology across animals}
\label{section-senses}

\textbf{Receptor types.}
Receptor cells are classified by the
stimulus energy they transduce and by
the transduction mechanism
\cite[Table 14.1]{hill2012animal}:
\begin{center}
\begin{tabular}{ll}
Receptor & Transduction \\
\hline
mechano-, vestibular,
& ionotropic \\
\quad auditory, osmo- & \\
thermo- & ionotropic \\
electro- & ionotropic \\
taste: salty, sour & ionotropic \\
taste: sweet, bitter, & metabotropic \\
\quad umami & \\
olfactory & metabotropic \\
 & \quad (vertebrates) \\
photo- & metabotropic \\
magneto- & unknown \\
\end{tabular}
\end{center}
The stimulus produces a graded receptor
potential, which in turn sets the
frequency of action potentials in the
afferent fibre. Which modality is felt
depends only on which receptor the fibre
comes from (M\"uller's law of specific
nerve energies)
\cite[Chapter 1]{schmidt1981sensory}.

\textbf{Adaptation.}
Under a sustained stimulus, \emph{tonic}
(slowly adapting) receptors keep firing,
while \emph{phasic} (rapidly adapting)
receptors fire mainly at the onset and
change of the stimulus. Merkel discs and
Ruffini endings are tonic; Meissner and
Pacinian corpuscles are phasic, the
Pacinian corpuscle extremely so because
of its lamellae \cite[Chapter
14]{hill2012animal}.

\textbf{Weber-Fechner law.}
The difference threshold $\Delta S$ (the
smallest detectable change of a stimulus
of size $S$) is roughly proportional to
$S$ (Weber's rule):
$$
\frac{\Delta S}{S} = k.
$$
For pressure on the skin $k$ is about
$3\%$. If Weber's rule holds, the
intensity of sensation $\psi$ grows with
the logarithm of the stimulus (Fechner):
$$
\psi = c \log\frac{S}{S_0},
$$
with $S_0$ the absolute threshold and $c$
a constant. This holds only over limited
ranges; Stevens' power function
$\psi = c\,(S - S_0)^n$ fits more
modalities \cite[Chapter
1]{schmidt1981sensory}. Receptor
potentials often grow with the logarithm
of intensity too, which compresses a huge
range of intensities into a usable
response \cite[Chapter
7]{randall1997eckert}.

\textbf{Smell.}
Olfactory qualities resist clean
classification; a conventional list of
primary odours is flowery, ethereal,
musky, camphorous, sweaty, rotten and
stinging \cite[Chapter
8]{schmidt1981sensory}.

\textbf{Eyes.}
\begin{itemize}
\item \emph{Compound eyes} (insects and
many crustaceans) are made of ommatidia,
each with a cuticular lens, a crystalline
cone and about eight retinular cells
whose microvilli form a rhabdom
containing rhodopsin. In
\emph{apposition} eyes (mostly diurnal
insects) each ommatidium is optically
isolated by pigment and receives light
from its own small angle, giving a
mosaic image. In \emph{superposition}
eyes (crepuscular and nocturnal insects,
deep-water crustaceans) light from one
direction reaches several rhabdoms:
optically, through a clear zone between
cones and rhabdoms, or neurally, by
convergence of receptors that look at the
same point. Superposition trades acuity
for sensitivity \cite[Section
9.8.4]{willmer2005environmental}.
\item \emph{Camera eyes} of cephalopods
and vertebrates evolved independently
but are strikingly alike: cornea, iris,
a lens, and a retina with a fovea. In
cephalopods the receptors point toward
the lens and the light; in vertebrates
light crosses the retinal neurons before
reaching the rods and cones, and the
axons leave through the optic disc,
causing a blind spot \cite[Section
9.8.4]{willmer2005environmental},
\cite[Chapter 14]{hill2012animal}.
\item Invertebrate (rhabdomeric)
photoreceptors depolarize in light,
whereas vertebrate rods and cones
hyperpolarize
\cite[Section
9.8.4]{willmer2005environmental}.
\end{itemize}

\textbf{Electroreception.}
Electroreceptors work only in water. They
are modified lateral-line hair cells.
\emph{Ampullary} organs (the ampullae of
Lorenzini of sharks and rays) sense weak,
low-frequency fields, for instance the
muscle potentials of prey buried in sand;
\emph{tuberous} organs of weakly electric
fish sense the high-frequency discharge
of the fish's own electric organ, used
for electrolocation and communication
(active electroreception)
\cite[Section
9.8.6]{willmer2005environmental},
\cite[Chapter 7]{randall1997eckert}.

\section{Organization of nervous systems}
\label{section-nervous-system}

\textbf{Invertebrate ganglia.}
An arthropod ganglion has an outer rind of
neuronal cell bodies, without synapses,
and an inner core with tracts of axons
and the \emph{neuropil}, the region of
synaptic contacts between axons and
dendrites \cite[Chapter
15]{hill2012animal}.

\textbf{Vertebrates.}
\begin{itemize}
\item The central nervous system is the
brain and spinal cord; the peripheral
nervous system is the nerves and
ganglia, with sensory and motor
divisions.
\item The somatic system controls
skeletal muscle. The autonomic system
controls internal effectors and has a
sympathetic (fight or flight), a
parasympathetic (rest and digest) and an
enteric division (gut, largely
independent of the brain)
\cite[Chapter 15]{hill2012animal}.
\item The brain has a forebrain
(telencephalon, diencephalon), a midbrain
(mesencephalon) and a hindbrain
(metencephalon, myelencephalon)
\cite[Chapter 15]{hill2012animal}. The
diencephalon contains the hypothalamus
(homeostasis, endocrine control, the
circadian clock); the hindbrain
contains the cerebellum (motor
coordination) and the medulla
(breathing, autonomic reflexes).
\item Glia: Schwann cells (peripheral)
and oligodendrocytes (central) make
myelin; astrocytes give metabolic
support, take up K$^+$ and
transmitters; microglia are phagocytes
related to immune cells
\cite[Chapter 12]{hill2012animal}.
\item The human brain has tens of
billions of neurons and up to five
times as many glial cells
\cite[Section
9.2.1]{willmer2005environmental}.
\end{itemize}

\section{Endocrine system}
\label{section-endocrine}

A hormone is a chemical secreted into the
blood that acts on receptors of distant
target cells. Hill sorts hormones into
three chemical classes
\cite[Chapter 16]{hill2012animal}:
\begin{itemize}
\item steroids, made from cholesterol
(gonads, adrenal cortex, ecdysone in
arthropods), lipid-soluble;
\item peptides and proteins (TRH,
insulin, growth hormone), water-soluble;
\item amines, modified amino acids:
melatonin (from tryptophan),
catecholamines and thyroid hormones
(from tyrosine); the thyroid hormones are
lipid-soluble, the others water-soluble.
\end{itemize}
Lipid-soluble hormones cross the membrane,
bind intracellular receptors and mostly
change gene transcription, so they act
slowly; they travel in the blood bound to
carrier proteins. Water-soluble hormones
bind membrane receptors and act through
second messengers, usually faster. Some
steroids also act on membrane receptors
\cite[Chapter 16]{hill2012animal}.
Juxtacrine signalling is between cells in
direct contact.

\textbf{Hypothalamus and pituitary.}
Hypothalamic neurosecretory cells release
releasing and inhibiting hormones into
the hypothalamo-hypophysial portal
system, which carries them to the
anterior pituitary
\cite[Chapter 16]{hill2012animal}.
Releasing hormones are CRH, TRH, GHRH
and GnRH; inhibiting hormones are
somatostatin (on GH), dopamine (on
prolactin), GnIH and MSH-IH. Ghrelin
also stimulates GH, and kisspeptins drive
the GnRH neurons
\cite[Appendix K]{hill2012animal}.
\begin{center}
\begin{tabular}{lll}
Hormone & Target & Main control \\
\hline
TSH & thyroid & TRH; thyroid \\
 & & \quad hormones \\
ACTH & adrenal cortex & CRH; \\
 & & \quad glucocorticoids \\
GH & bone, muscle, & GHRH, ghrelin; \\
 & \quad fat, liver & \quad somatostatin \\
prolactin & mammary gland & dopamine \\
FSH & gonads & GnRH \\
LH & gonads & GnRH \\
MSH & melanocytes & MSH-IH \\
\end{tabular}
\end{center}
FSH drives sperm production and follicle
growth; LH drives sex-steroid production,
ovulation and the corpus luteum. Thyroid
hormones and glucocorticoids feed back
negatively on their own axes
\cite[Chapter 16, Appendix
K]{hill2012animal}.

\section{Circadian rhythms and melatonin}
\label{section-clocks}

A circadian rhythm has a period of about a
day and is endogenous: in constant
conditions it free-runs, and a light-dark
cycle entrains it to exactly $24$ hours.
In mammals the master clock is the pair
of suprachiasmatic nuclei (SCN) of the
hypothalamus, which receive light
information from the retina
\cite[Chapter 15]{hill2012animal}.

The SCN controls the pineal gland, which
makes melatonin from tryptophan and
secretes it at night in diurnal and
nocturnal mammals alike, so melatonin is
called the hormone of darkness
\cite[Chapter 15]{hill2012animal}. The
length of nightly secretion encodes the
photoperiod, and melatonin times seasonal
reproduction through its effects on the
GnRH-LH-FSH axis \cite[Chapter
17]{hill2012animal}.

\section{The practicals}
\label{section-practicals}

\textbf{Cricket electroretinogram.}
One electrode was placed on the compound
eye of the house cricket \emph{Acheta
domesticus} and a reference electrode on
the abdomen, inside a Faraday cage. The
eye was stimulated with white light at
$25$, $50$, $75$ and $100\%$ of a
calibrated lamp, then through blue and
red cellophane filters, four trials each,
with a final white series as a check for
deterioration. ANOVA across intensities
was significant for white
($p = 0.0007$) and red ($p < 0.001$)
light but not for blue ($p = 0.61$).
Blue light gave the largest responses at
all intensities, which the report reads
as saturation of the eye at short
wavelengths.

\textbf{Sensory tests.}
On fifteen people the class measured, at
each phalanx of each finger, on the palm
and on the forearm, a touch
discrimination value (probably the
two-point threshold in centimetres; the
files do not give the unit), at room
temperature and after cooling. Palm
($\approx 0.75$) and forearm
($\approx 1.0$) differed significantly
from the fingers ($\approx 0.25$ to
$0.5$) in Tukey comparisons; cooling
raised the values slightly, mostly not
significantly. Other tests timed
repeated tasks over six trials (times fell
from about $50$ to $40\ \mathrm{s}$, a
learning effect) and measured a ruler
drop in centimetres, a reaction-time test.

\textbf{Cardiovascular practical.}
Students recorded an electrocardiogram at
rest, during five minutes of exercise and
in recovery, together with blood pressure
and pulse. Resting heart rates of men and
women did not differ significantly in a
two-tailed unpaired $t$ test; the class
discussed that a one-tailed test, halving
the $p$-value, is justified only when
theory predicts the direction. Rhythm was
displayed by a Poincar\'e return map,
plotting each beat period $P_n$ against
the next one $P_{n+1}$. During exercise
the cloud of points is smaller and closer
to the origin: the periods are shorter
and less variable. The coefficient of
variation of the periods gives a number
for this spread.

\section{Water, salts and excretion}
\label{section-osmoregulation}

This section is written from the
textbooks; the surviving coursework has no
material on it.

\textbf{Regulation and conformity.}
Osmoconformers let their blood osmotic
pressure follow the environment;
osmoregulators keep it nearly constant.
Animals tolerating a narrow or wide range
of salinity are stenohaline or euryhaline
\cite[Chapter 27]{hill2012animal}.

\textbf{Environments.}
\begin{itemize}
\item All freshwater animals are
hyperosmotic regulators: they gain water
osmotically, produce dilute urine and
take up ions actively.
\item Most marine invertebrates and the
hagfish are close to isosmotic with
seawater. Marine teleosts are hyposmotic
regulators: they lose water, drink
seawater and excrete salt through the
gills.
\item Migratory fish such as salmon
switch between the two modes
(hyper-hyposmotic regulation).
\end{itemize}
See \cite[Chapter 28]{hill2012animal}.

\textbf{Kidneys and nitrogen.}
Urine forms by filtration (or secretion)
of a primary urine followed by
reabsorption and secretion along the
tubule; in insects the Malpighian tubules
secrete it. Nitrogen leaves as ammonia
(ammonotelic, mostly water breathers; toxic
but cheap), urea (ureotelic, mammals;
less toxic, costs ATP) or uric acid
(uricotelic, birds, reptiles, insects;
insoluble, saves water)
\cite[Chapter 29]{hill2012animal}.

\bibliography{my}
\bibliographystyle{amsalpha}

\end{document}
